
\(f_\psi\) de \(K[X_i]_{i\geq 1}\) dans \(A\) tel que \(\Ghat(f_\psi)(\phi)=\psi\).

1.3. On définit pour tout \(A\in\operatorname{Ob}(\mathcal C)\) une filtration sur \(\Ghat(A)\) de la façon suivante : une série \(\psi=1+a_1t+\ldots\) sera dite de filtration \(\geq k\) si \(a_i=0\) pour \(1\leq i<k\). On note \(F^k\Ghat(A)\) le sous-groupe ainsi obtenu ; c'est un foncteur sur \(\mathcal C\), et nous considérerons désormais \(\Ghat\) comme foncteur à valeurs dans la catégorie \(\underline{\mathrm{Ab}}.\underline{\mathrm{fil}}.\) des groupes abéliens filtrés grâce à cette filtration(avec \(F^0\Ghat=\Ghat\)).

Soit \(H\) un foncteur sur \(\mathcal C\) à valeurs dans \(\underline{\mathrm{Ab}}\,\underline{\mathrm{fil}}\)~; nous supposerons que \(H\) satisfait les deux conditions suivantes :
\[
\begin{array}{ll}
\text{i)} & \forall A\in\operatorname{Ob}(\mathcal C),\ H(A)\text{ est séparé et complet}.\\[0.3em]
\text{ii)} & \text{Si }A\in\operatorname{Ob}(\mathcal C)\text{ est une }K\text{-algèbre graduée, munie de la filtration}\\
& \text{associée à sa graduation, et si }G\text{ est un groupe fini opérant sur }A\\
& \text{par morphismes gradués, alors }H(A^G)\xrightarrow{\sim} H(A)^G\\
& \text{(l'exposant }G\text{ désignant les invariants sous }G\text{)}
\end{array}
\]

En fait, nous n'aurons besoin de ces conditions que lorsque \(A\) est un anneau de polynômes, et \(G\) un produit fini de groupes symétriques opérant par permutation des indéterminées.

\textbf{Exemples}~: 1)~\(H(A)=\Ghat(A)\). 2)~\(H(A)=1+A[[t]]^+\)~; ce foncteur sera noté \(\Gg\) (\(g\), initiale de « grossière »). 3)~\(H(A)=\Ga(A)\), où \(\Ga(A)\) désigne le groupe (pour l'addition) des séries formelles telles que \(a_0=0\), \(a_i\in\Fil^i A\) pour \(i\geq 1\), muni de la filtration analogue à celle de \(\Ghat\).

1.4. Pour tout \(n\), nous noterons \(f_n\) l'endomorphisme de \(K[X_1,\ldots,X_n]\)
